\chapter{Cadre général du projet}

\section{Introduction}

La sécurité applicative constitue aujourd'hui l'un des enjeux majeurs du développement
logiciel moderne. L'accélération des cycles de livraison, portée par l'adoption massive
des pratiques DevOps et des architectures microservices, a profondément modifié le
rapport des équipes techniques à la sécurité : celle-ci ne peut plus être traitée comme
une phase terminale du cycle de vie logiciel, mais doit être intégrée dès la conception
et automatisée à chaque étape du pipeline de déploiement.

C'est dans cette perspective que s'inscrit \textbf{InvisiThreat}, une plateforme
DevSecOps dédiée à la détection automatisée de vulnérabilités. Le projet vise à rendre
la sécurité applicative \textit{continue}, \textit{traçable} et \textit{actionnable},
en offrant aux équipes techniques un point d'entrée unique pour orchestrer des analyses
statiques et dynamiques, consolider les résultats et produire des rapports exploitables.

La détection repose sur deux familles d'analyses complémentaires. La
\textbf{SAST} (\textit{Static Application Security Testing}) inspecte le code source
afin d'identifier des patterns dangereux, des appels non sécurisés ou des configurations
à risque, sans nécessiter l'exécution de l'application. La
\textbf{DAST} (\textit{Dynamic Application Security Testing}), quant à elle, teste
l'application en cours d'exécution pour détecter les failles d'exposition réelles :
injections, mauvaises configurations HTTP, contournements d'authentification. Ces deux
approches sont complémentaires : la SAST couvre les défauts logiques en amont, tandis
que la DAST valide l'impact réel sur l'application déployée.

InvisiThreat unifie ces deux dimensions dans une architecture cohérente : backend
\textbf{FastAPI} avec \textbf{SQLAlchemy}, base de données \textbf{PostgreSQL}, moteurs
SAST via \textbf{Semgrep} et \textbf{Bandit}, moteur DAST via \textbf{OWASP ZAP} en
mode daemon, et frontend \textbf{React} enrichi de notifications temps réel via
\textbf{Socket.IO}.

Ce premier chapitre présente le cadre général du projet : organisme d'accueil, contexte
et problématique de sécurité, étude critique des solutions existantes, solution proposée
et méthodologie de développement adoptée.

% ─────────────────────────────────────────────────────────────────
\section{Présentation de l'organisme d'accueil}

\textbf{Mobelite} est une entreprise spécialisée dans la conception et le développement
d'applications web et mobiles à forte valeur ajoutée. Opérant dans un environnement
multi-projets à livraison rapide, Mobelite accompagne ses clients dans la transformation
numérique de leurs processus métiers, en délivrant des solutions sur mesure reposant sur
des technologies modernes et des architectures orientées services.

L'environnement technique de l'entreprise se caractérise par une adoption significative
des outils DevOps : gestion de versions sous \textbf{Git}, conteneurisation via
\textbf{Docker}, déploiements automatisés par pipelines CI/CD, et développement
d'applications exposant des \textbf{APIs REST} consommées par des frontends React ou
des clients mobiles. Cette stack technologique moderne, bien qu'efficace, génère une
surface d'attaque étendue : endpoints publics, gestion des tokens d'authentification,
intégration de bibliothèques tierces et exposition de services.

Avant la mise en place d'InvisiThreat, les audits de sécurité étaient réalisés de
manière ponctuelle et hétérogène, sans outillage unifié ni traçabilité structurée des
résultats. L'absence d'un historique centralisé des vulnérabilités par projet rendait
difficile la mesure de l'évolution du risque et la justification des efforts de
remédiation. Ce constat a motivé la création d'un outil interne capable de
\textbf{standardiser}, \textbf{automatiser} et \textbf{centraliser} la détection des
failles de sécurité à l'échelle de l'ensemble du portefeuille applicatif.

InvisiThreat a ainsi été conçu comme une réponse opérationnelle directe à ces besoins,
s'intégrant naturellement dans l'écosystème technique existant tout en apportant une
couche de sécurité continue, mesurable et actionnable.

% ─────────────────────────────────────────────────────────────────
\section{Présentation du sujet}

Le sujet de ce projet de fin d'études porte sur la conception, le développement et le
déploiement d'une plateforme DevSecOps unifiée, baptisée \textbf{InvisiThreat}, capable
d'orchestrer des analyses de sécurité SAST et DAST, de normaliser et consolider les
résultats issus de moteurs hétérogènes, et de produire des rapports exploitables par
les équipes techniques et managériales.

La plateforme s'appuie exclusivement sur des outils open source reconnus dans l'industrie
de la sécurité applicative. L'analyse statique est assurée par \textbf{Semgrep}, moteur
de règles polyglotte à haute précision, et \textbf{Bandit}, analyseur dédié aux sources
Python. L'analyse dynamique est orchestrée via l'\textbf{API REST d'OWASP ZAP} opérant
en mode daemon, permettant l'automatisation complète des phases de spider et de scan
actif sans intervention manuelle.

L'ambition centrale est de proposer un \textbf{point d'entrée unique} pour le lancement,
le suivi et la priorisation des vulnérabilités, en supprimant la fragmentation des
rapports entre outils et en fournissant une vue consolidée du risque par projet, par
version et dans le temps.

\subsection{Contexte du projet}

La croissance du nombre de vulnérabilités publiées illustre l'ampleur du défi auquel
font face les équipes de sécurité. En 2023, la base nationale de vulnérabilités NVD
(\textit{National Vulnerability Database}) a enregistré \textbf{29\,065 CVE} publiées,
un record historique à cette date \cite{nvd2023}. Cette tendance se confirme en 2024,
avec plus de \textbf{40\,000 CVE} anticipées selon les projections du NIST, traduisant
une augmentation d'environ 38\,\% par rapport à l'année précédente \cite{nvd2024}.

Face à cette pression, les organisations adoptant des pratiques CI/CD sont
particulièrement exposées : chaque déploiement peut introduire de nouvelles dépendances
vulnérables ou modifier la surface d'attaque d'une API. L'\textbf{OWASP API Security
Top 10} (édition 2023) identifie les risques les plus critiques pesant sur les APIs
REST modernes, parmi lesquels les défauts d'autorisation au niveau des objets
(\textit{BOLA}), les mauvaises configurations d'authentification et l'exposition
excessive de données \cite{owaspApi2023}. Ces catégories sont directement applicables
aux applications développées dans l'environnement Mobelite.

Le contexte applicatif d'InvisiThreat combine des stacks variées : APIs REST FastAPI,
frontends React, services tiers intégrés via webhooks et tokens. Cette diversité
technologique augmente la surface d'attaque et rend indispensable une automatisation
des audits capable de couvrir à la fois la qualité du code source et le comportement
de l'application déployée.

\subsection{Problématique}

Les pratiques de sécurité traditionnelles reposent sur des audits manuels périodiques,
réalisés par des équipes spécialisées intervenant ponctuellement sur le périmètre
applicatif. Cette approche, bien qu'efficace dans un contexte de livraisons peu
fréquentes, devient structurellement inadaptée face aux cadences imposées par les
pipelines CI/CD modernes, où plusieurs déploiements peuvent intervenir quotidiennement.

La multiplication des outils de scan isolés aggrave ce problème. Chaque moteur génère
ses propres rapports, dans ses propres formats, avec ses propres nomenclatures de
sévérité. L'analyste sécurité doit alors \textbf{corréler manuellement} des résultats
hétérogènes, sans garantie d'exhaustivité ni de cohérence entre les versions analysées.

Les limitations intrinsèques de chaque approche renforcent cette problématique :

\begin{itemize}
    \item La \textbf{SAST seule} est insuffisante : elle identifie des vulnérabilités
    potentielles au niveau du code source, mais ne peut pas confirmer leur exploitabilité
    réelle en production, ni détecter les failles de configuration de l'environnement
    d'exécution.

    \item La \textbf{DAST seule} ne détecte pas les vulnérabilités logiques enfouies
    dans le code, notamment les injections de second ordre, les défauts de contrôle
    d'accès ou les mauvaises pratiques cryptographiques. Sa couverture dépend fortement
    de la qualité de la navigation et des règles d'attaque paramétrées.

    \item L'absence de \textbf{traçabilité historique} par projet empêche de mesurer
    objectivement la réduction du risque au fil des correctifs appliqués, rendant
    difficile la démonstration de la maturité sécurité à des parties prenantes externes.
\end{itemize}

La problématique centrale du projet est donc la suivante : \textbf{comment orchestrer
de manière cohérente et automatisée des analyses SAST et DAST, normaliser leurs
résultats selon une nomenclature unifiée et fournir aux équipes techniques une vision
priorisée et historisée du risque applicatif, intégrable dans un pipeline DevOps
existant~?}

\subsection{Étude de l'existant}

Afin de positionner InvisiThreat par rapport à l'écosystème existant, une analyse
comparative des outils de sécurité applicative les plus répandus a été conduite.
Ces outils se répartissent en trois catégories : SAST, DAST et SCA
(\textit{Software Composition Analysis}). Le tableau~\ref{tab:tool-comparison} en
présente une synthèse selon plusieurs critères d'évaluation.

\begin{table}[H]
\centering
\caption{Comparatif des outils de scan de sécurité applicative}
\label{tab:tool-comparison}
\renewcommand{\arraystretch}{1.4}
\begin{tabularx}{\textwidth}{>{\bfseries}l l c l l X}
\hline
\rowcolor{gray!15}
Outil & Type & Open Source & Intégration & Reporting & Limitations principales \\
\hline
SonarQube   & SAST       & Partiel & CI/CD, API  & HTML, JSON        & Absence de DAST ; version Community limitée ; coûts Enterprise élevés \\
Semgrep     & SAST       & Oui     & CLI, CI/CD  & SARIF, JSON       & Faux positifs si les règles ne sont pas adaptées au contexte \\
Bandit      & SAST       & Oui     & CLI         & JSON, TXT         & Limité aux sources Python ; pas de support multi-langages \\
OWASP ZAP   & DAST       & Oui     & API, CLI    & HTML, JSON, XML   & Paramétrage complexe ; faux négatifs sur SPAs ou authentifications avancées \\
Burp Suite  & DAST       & Partiel & GUI, API    & HTML, XML         & Licences élevées ; automatisation complète réservée à l'édition Pro \\
Snyk        & SCA / SAST & Partiel & CI/CD, API  & HTML, JSON        & Couverture SAST tributaire du plan tarifaire \\
Checkmarx   & SAST       & Non     & CI/CD, API  & PDF, HTML         & Coût prohibitif ; modèle boîte noire ; déploiement complexe \\
\hline
\end{tabularx}
\end{table}

L'analyse de ce panorama met en évidence une lacune commune à l'ensemble de ces
solutions : aucune ne couvre simultanément l'analyse statique, l'analyse dynamique
et la restitution unifiée dans une interface unique accessible sans infrastructure
dédiée coûteuse.

\subsection{Critique de l'existant}

Au-delà des limitations individuelles de chaque outil, c'est leur utilisation conjointe
qui révèle les principales faiblesses des approches actuelles. Les rapports générés par
des moteurs différents utilisent des formats hétérogènes — SARIF, JSON propriétaire,
HTML — et des nomenclatures de sévérité incompatibles, rendant toute comparaison
directe laborieuse et sujette à interprétation.

La corrélation entre une vulnérabilité détectée par SAST et sa confirmation en DAST
n'est pratiquement jamais automatisée dans les outils du marché. Cette absence de
corrélation oblige les analystes à traiter les résultats en silos, ce qui allonge
les délais de remédiation et introduit un risque de doublons ou d'omissions.

Par ailleurs, l'absence d'un \textbf{historique centralisé} par projet constitue un
frein à la mesure de la maturité sécurité dans le temps. Les solutions SaaS comme
Snyk ou Checkmarx proposent certes des tableaux de bord, mais au prix d'une dépendance
à des infrastructures cloud externes et de coûts de licence significatifs, incompatibles
avec une politique de maîtrise des données et des budgets opérationnels contraints.

Enfin, la complexité de paramétrage des outils DAST — en particulier pour les
applications à authentification avancée ou les \textit{Single Page Applications} — est
rarement documentée de façon opérationnelle, ce qui limite leur adoption et leur
efficacité en environnement de production.

\subsection{Solution proposée}

Face aux lacunes identifiées, \textbf{InvisiThreat} propose une approche intégrée
reposant sur quatre principes fondateurs :

\begin{enumerate}
    \item \textbf{Pipeline unifié} : un point d'entrée unique pour déclencher des scans
    SAST (Semgrep, Bandit) et DAST (OWASP ZAP), sans multiplier les interfaces ni les
    configurations manuelles.

    \item \textbf{Normalisation des résultats} : toutes les alertes, quelle que soit
    leur origine, sont converties vers un modèle interne commun avec une sévérité
    normalisée sur cinq niveaux (\textit{Critical}, \textit{High}, \textit{Medium},
    \textit{Low}, \textit{Informational}), accompagnées de preuves techniques
    (extraits de code, requêtes/réponses HTTP).

    \item \textbf{Traçabilité historique} : chaque exécution est rattachée à un projet,
    une version applicative, un outil et un horodatage précis, permettant des
    comparaisons entre scans successifs et la mesure objective de la réduction du risque.

    \item \textbf{Accessibilité opérationnelle} : notifications temps réel via
    Socket.IO, export PDF et JSON des rapports, et exposition d'une API REST documentée
    permettant l'intégration dans des pipelines CI/CD existants.
\end{enumerate}

Ce positionnement vise à rendre la sécurité applicative actionnable pour les
développeurs et les chefs de projet, sans dépendance à des solutions propriétaires
ou à des infrastructures cloud tierces.

% ─────────────────────────────────────────────────────────────────
\section{Méthodologie de développement}

\subsection{Introduction}

Le développement d'InvisiThreat implique l'orchestration de composants techniques
hétérogènes — moteurs de scan, API, base de données, frontend réactif — dont les
dépendances et les comportements ne sont pas entièrement prévisibles dès le départ.
Une méthodologie capable d'absorber les incertitudes techniques et d'intégrer
progressivement les retours des utilisateurs s'est naturellement imposée.

\subsection{Méthodologie Agile}

La méthode \textbf{Agile} a été retenue pour sa capacité à structurer le développement
de manière incrémentale, en validant chaque fonctionnalité avant de progresser vers la
suivante. Elle offre une réponse adaptée à la nature exploratoire d'un projet comme
InvisiThreat, où les paramètres de scan, les règles de détection et les besoins métiers
évoluent au fil des itérations.

Le \textbf{backlog produit} a été structuré en cinq épics principaux :
\textit{authentification et gestion des accès}, \textit{gestion des projets et cibles},
\textit{orchestration des scans SAST}, \textit{orchestration des scans DAST}, et
\textit{reporting et notifications}. Chaque \textit{user story} est associée à des
critères d'acceptation mesurables, par exemple : \og lancer un scan DAST sur une URL
cible et obtenir un rapport complet en moins de cinq minutes \fg{}, ou \og recevoir
une notification Socket.IO en temps réel à la détection d'une vulnérabilité critique \fg{}.

\subsection{L'approche Scrum}

Parmi les cadres Agile disponibles, \textbf{Scrum} a été choisi pour sa structure
itérative claire et la visibilité qu'il offre sur les livrables à chaque incrément.
Sa cadence régulière de \textit{sprints} a permis de prioriser les fonctionnalités
à fort risque technique — notamment l'intégration de l'API ZAP en mode daemon et la
normalisation des vulnérabilités issues de moteurs hétérogènes — avant d'aborder les
couches de présentation et d'export.

Le processus Scrum adopté s'articule autour des cérémonies standard : planification de
sprint, réunions quotidiennes (\textit{daily stand-up}), revue de sprint et
rétrospective. Les rôles sont distribués entre le \textit{Product Owner} — responsable
de la priorisation du backlog selon la valeur métier et le risque —, le
\textit{Scrum Master} — garant du respect du cadre et de la levée des obstacles — et
l'\textit{équipe de développement}, en charge de la réalisation des incréments.

La figure~\ref{fig:scrum-diagram} illustre le cycle Scrum tel qu'il a été appliqué
au projet InvisiThreat, depuis la construction du \textit{Product Backlog} jusqu'à
la livraison des incréments lors des revues de sprint.

\begin{figure}[H]
\centering
\fbox{\rule{0pt}{2.5in}\rule{0.9\linewidth}{0pt}}
\caption{Cycle Scrum appliqué au projet InvisiThreat.}
\label{fig:scrum-diagram}
\end{figure}

\noindent Le code PlantUML correspondant au diagramme Scrum est fourni ci-dessous
à titre de référence pour la génération du schéma :

\begin{verbatim}
@startuml
skinparam backgroundColor #FAFAFA
skinparam rectangleBorderColor #2C3E50
skinparam arrowColor #2980B9

actor "Product Owner" as PO
actor "Scrum Master"  as SM
actor "Équipe Dev"    as DEV

rectangle "Product Backlog\n(Épics + User Stories)" as PB #LightBlue
rectangle "Sprint Backlog\n(Tâches du sprint)"      as SB #LightGreen
rectangle "Incrément livrable\n(Fonctionnalité testée)" as INC #LightYellow

PO  --> PB  : Prioriser et affiner
SM  --> SB  : Planifier le sprint (2 semaines)
DEV --> SB  : Réaliser et tester
SB  --> INC : Livrer l'incrément
INC --> PO  : Revue et validation
INC --> PB  : Rétroaction et ajustements
@enduml
\end{verbatim}

\subsection{Planification des sprints}

Le projet a été structuré en \textbf{cinq sprints} de deux semaines chacun, couvrant
l'ensemble du périmètre fonctionnel d'InvisiThreat. Le tableau~\ref{tab:sprints}
détaille les objectifs, le contenu et les livrables concrets de chaque itération.

\begin{table}[H]
\centering
\caption{Planification et livrables des sprints InvisiThreat}
\label{tab:sprints}
\renewcommand{\arraystretch}{1.5}
\begin{tabularx}{\textwidth}{c >{\bfseries}l X X}
\hline
\rowcolor{gray!15}
\textbf{Sprint} & \textbf{Objectif principal} & \textbf{Contenu} & \textbf{Livrables} \\
\hline
1 &
Socle technique &
Configuration Docker Compose (4 services), modèle de données initial, migrations Alembic, structure des répertoires backend et frontend. &
Environnement conteneurisé opérationnel, schéma PostgreSQL initial, pipeline CI vide GitHub Actions. \\

\hline
2 &
Authentification et gestion des projets &
Inscription, connexion JWT, OTP via QR code (TOTP), gestion des rôles RBAC, CRUD projets et URLs cibles. &
Module d'authentification sécurisé, API projets fonctionnelle, interface de connexion React. \\

\hline
3 &
Intégration SAST &
Exécution Semgrep et Bandit via sous-processus Python, parsing JSON des résultats, normalisation et insertion des vulnérabilités en base. &
Pipeline SAST opérationnel, modèle \texttt{Vulnerability} peuplé, endpoint \texttt{POST /api/scans/sast}. \\

\hline
4 &
Intégration DAST &
Connexion à l'API OWASP ZAP en mode daemon, spider automatique, scan actif, récupération des alertes, stockage des preuves HTTP. &
Pipeline DAST opérationnel, alertes ZAP normalisées, endpoint \texttt{POST /api/scans/dast}. \\

\hline
5 &
Dashboard, reporting et notifications &
Tableau de bord React par sévérité, export PDF/JSON, notifications Socket.IO (fin de scan, vulnérabilité critique), emails transactionnels. &
Interface complète, export rapport fonctionnel, notifications temps réel validées. \\
\hline
\end{tabularx}
\end{table}

\subsection{Conclusion}

L'organisation en cinq sprints progressifs a permis de construire InvisiThreat de
manière maîtrisée, en sécurisant d'abord les fondations techniques avant d'intégrer
les modules de scan puis les couches de présentation et de notification. La cadence
Scrum a favorisé une validation régulière des fonctionnalités critiques et une
réactivité face aux ajustements techniques imposés par l'intégration des moteurs de
scan. Ce cadre méthodologique a directement contribué à la qualité et à la robustesse
de la solution livrée.